\begin{resumo}[Abstract]
 \begin{otherlanguage*}{english}
   The v-FES cycling project, carried out by the BEM-TE-VI Technology Platform of the University of Brasília,
   follows its participants through regular assessments performed by a team of physiotherapists who use seven
   instruments. The Bem-Te-Vi app gathers these forms, but handles each of them separately, offers assessment
   history and comparison only in the ICF flow and has problems in data recording. This work defines the
   requirements and an interoperable data model in HL7 FHIR to structure these data and support the follow-up of
   assessments, taking as its starting point the ICF module implemented by the author. The research follows the
   design science method proposed by Wieringa and draws on the app's source code, the project documentation and
   the specialized literature on these instruments, without consulting the physiotherapy team. Sixty-seven
   requirements were specified, and the model, built in FHIR R4, represented all ICF structures, with instances
   of fictitious data approved by the validator. The analysis also showed that the app already automates the
   calculations within each assessment, although some of them deviate from the instrument, but does not automate
   the follow-up across assessments in the other forms. These results guide the implementation of the database
   and the validation with the physiotherapy team in the second part of this work.

   \vspace{\onelineskip}

   \noindent
   \textbf{Keywords}: Requirements engineering. Clinical data. ICF. HL7 FHIR. Physiotherapy.
 \end{otherlanguage*}
\end{resumo}
